% ------------------------------------------------------------------
% E4 (학습 중 주입) — 논문 본문에서 뺀 내용 보관 파일
%
% 2026-11-06 main.tex에서 분리. 다음 논문/보고서에서 쓰기 위해 원문을 그대로 둔다.
% 각 조각 앞의 주석은 main.tex의 어느 자리에 있었는지를 적은 것이다.
%
% 데이터 : src/results/single_mnist_linear_train_*.csv, single_mnist_cnn_train_*.csv
% 실행   : python src/run_e1_models.py --models mnist_linear --mode train \
%            --p-list 0 1e-4 3e-4 1e-3 3e-3 1e-2 3e-2 1e-1 --defenses none clip parity tmr --trials 10
%          python src/run_e1_models.py --models mnist_linear --mode train --rad leo_saa \
%            --train-hours 24 168 720 2160 --defenses none clip --trials 10
%          python src/run_e1_models.py --models mnist_cnn --mode train --p-list 0 1e-3 1e-2 --defenses none clip --trials 5
% 그림·표 : paper/make_figures.py 의 load_e4 / fig_e4 / table_e4 가 지금도
%          figures/fig_e4_train.pdf, tables/table_e4.tex, tables/table_e4_leo.tex,
%          그리고 numbers.tex 의 \Efour... 매크로를 만든다 (본문에서 안 쓸 뿐).
% 되살리기: 아래 조각을 주석에 적힌 자리에 붙여 넣으면 된다. \ref{sec:e4}, \ref{fig:e4},
%          \ref{tab:e4}, \ref{tab:e4leo} 라벨은 4.9절 조각 안에 있다.
% ------------------------------------------------------------------

%==================================================================
% 위치: 초록 (마지막 문장)
%==================================================================
학습 도중에 주입하면(E4) 치명 확률은 같지만 실패가 전부 NaN 붕괴로 나타나 방어 없는 궤도상 학습은 수 일 안에 반드시 죽는 반면, 클리핑을 켜면 SGD가 잘린 가중치를 다시 고쳐 학습이 추론보다 오히려 강건해졌다.

%==================================================================
% 위치: 3.4 실험 설계 — E4 문단 (E3 문단 뒤, E5 문단 앞)
%==================================================================
\textbf{E4 (학습 중 주입).} 추론이 아니라 SGD 학습 도중에 명중을 넣는다. stress 모드에서는 총 $pN_w$개의 사건을 전체 최적화 스텝에 고르게 분산해 넣고(수업 코드처럼 epoch마다 주는 것이 아니라 전체 학습에 한 번의 선량), \code{leo\_saa} 모드에서는 학습 전체를 $T$일의 궤도 노출로 가정해 사건 시각을 스텝에 대응시키므로 SAA 통과 구간의 명중이 특정 스텝 구간에 몰린다. 방어는 매 주입 직후에 적용한다. \code{mnist\_linear}(60{,}000장, 2 epoch)에 stress $p\in\{0,10^{-4},3\times10^{-4},10^{-3},3\times10^{-3},10^{-2},3\times10^{-2},10^{-1}\}$, 방어 4종(없음, 클리핑, 패리티, TMR), 각 10회, 그리고 \code{leo\_saa} $T\in\{1,7,30,90\}$일, 방어 2종, 각 10회를 돌렸다. \code{mnist\_cnn}은 $p\in\{0,10^{-3},10^{-2}\}$, 방어 2종, 각 5회이다. 기준 정확도는 $p=0$(방사선 없음)으로 학습한 10회의 중앙값이고, 학습 도중 손실·가중치·그래디언트에 NaN/$\infty$가 생기면 붕괴로 기록하고 학습을 멈춘다.

%==================================================================
% 위치: 4.7 궤도 모드 (E3) 끝 문장
%==================================================================
학습 중 주입(E4, \ref{sec:e4}절)에서도 생존률은 같은 모형으로 예측되었다.

%==================================================================
% 위치: 4.9 절 전체 (4.8 방어 비교 E5 뒤, 4.10 비트별 치명도 E2 앞)
%==================================================================
\subsection{학습 중 주입 (E4): 전부 아니면 전무}
\label{sec:e4}
그림~\ref{fig:e4}와 표~\ref{tab:e4}, \ref{tab:e4leo}는 학습 도중에 명중을 넣은 결과다. 방사선 없이 학습한 기준 정확도는 \EfourCtrl{}\,\%였다. 방어가 없을 때 생존률이 50\,\%를 가로지르는 명중 수는 약 \EfourKfiftyNone{}개로 추론 중(모델 B·A 적합값 \KfiftyBNoneLinear{}--\KfiftyALinear{}개)과 같은 자릿수였고, 궤도 모드의 생존률(표~\ref{tab:e4leo})도 추론용 모델 A의 예측과 맞았다. 즉 ``명중 하나가 치명적일 확률''은 학습 중에도 비슷하다. 그러나 \emph{실패의 모양}이 다르다. 추론 중에는 실패한 시행이 10--80\,\%에 넓게 퍼졌지만(그림~\ref{fig:bimodal}), 학습 중에는 실패가 거의 전부 NaN 붕괴였다. $p=0.1\,\%$에서 중앙값 \EfourOneEmThreeNoneMed{}\,\%, 생존률 \EfourOneEmThreeNoneSurv{}, 붕괴율 \EfourOneEmThreeNoneCrash{}이고, $p\ge1\,\%$에서는 10회 모두 붕괴하여 정확도가 찍기 수준 10\,\%였다. 추론에서는 $10^{38}$ 크기의 가중치가 그 경로를 지나는 입력만 망치지만, 학습에서는 그 가중치가 만든 거대 로짓이 손실을 $\infty$로, 그래디언트를 NaN으로 만들고 다음 스텝의 옵티마이저가 \emph{모든} 가중치를 NaN으로 덮어쓴다. 수업 11주차에서 학생들이 관찰한 ``즉사''가 이것이다.

방어가 있으면 반대로 학습이 추론보다 \emph{더} 강건했다. 클리핑은 $p=10\,\%$(명중 약 10{,}000개)에서도 생존률 1.0, 정확도 \EfourOneEmOneClipMed{}\,\%로, 같은 선량의 추론(\pOneEmOneClipMed{}\,\%)보다 훨씬 좋았다. $\pm3$으로 잘린 가중치가 추론에서는 그대로 남아 누적 손상이 되지만, 학습에서는 이어지는 SGD 스텝이 그 가중치를 다시 고치기 때문이다. 패리티와 TMR은 모든 선량에서 기준 정확도 그대로였다. \code{mnist\_cnn}에서도 같았다. 방어 없음은 $p=0.1\,\%$(명중 약 220개)에서 5회 모두 붕괴했고, 클리핑은 $p=1\,\%$에서도 \EfourCnnOneEmTwoClipMed{}\,\%(기준 \EfourCnnCtrl{}\,\%)였다.

E4의 결론은 두 가지다. 첫째, 궤도에서 학습하는 신경망은 방어 없이는 수 일 안에 반드시 죽고, 죽을 때는 서서히가 아니라 한 번에 죽는다. 둘째, 그러나 가장 단순한 클리핑만 켜면 학습은 추론보다 오히려 안전하다. ``지상에서 학습하고 궤도에서는 추론만''이라는 통념은 방어 없는 경우에만 옳고, 클리핑이 있다면 궤도상 재학습(on-orbit fine-tuning)도 가능하다는 뜻이다.

\begin{figure}[!htb]
  \centering
  \includegraphics[width=0.95\linewidth]{fig_e4_train}
  \caption{E4: 학습 중 주입. (a) \code{mnist\_linear}에서 방어 4종의 학습 후 정확도. (b) 같은 명중 수에서 학습 중(실선)과 추론 중(점선)의 생존률 비교.}
  \label{fig:e4}
\end{figure}

\begin{table}[!htb]
  \centering\scriptsize
  \caption{E4 (\code{mnist\_linear}, 각 10회): 학습 후 중앙값 정확도(괄호는 생존률). 기준 정확도 \EfourCtrl{}\,\%. $\bar{k}$는 학습 전체에 걸친 기대 명중 수.}
  \label{tab:e4}
  \setlength{\tabcolsep}{3pt}
\begin{tabular}{lrrrrrrr}
\toprule
방어 & $p$=0.01\% & $p$=0.03\% & $p$=0.1\% & $p$=0.3\% & $p$=1\% & $p$=3\% & $p$=10\% \\
 & $\bar{k}\approx$10 & $\bar{k}\approx$30 & $\bar{k}\approx$102 & $\bar{k}\approx$305 & $\bar{k}\approx$1,016 & $\bar{k}\approx$3,049 & $\bar{k}\approx$10,163 \\
\midrule
없음 & 96.4 (1.0) & 96.2 (0.7) & 95.2 (0.6) & 10.0 (0.2) & 10.0 (0.0) & 10.0 (0.0) & 10.0 (0.0) \\
클리핑 & 96.4 (1.0) & 96.5 (1.0) & 96.4 (1.0) & 96.4 (1.0) & 96.4 (1.0) & 96.1 (1.0) & 95.3 (1.0) \\
패리티 & 96.4 (1.0) & 96.4 (1.0) & 96.4 (1.0) & 96.4 (1.0) & 96.4 (1.0) & 96.4 (1.0) & 96.4 (1.0) \\
TMR & 96.4 (1.0) & 96.4 (1.0) & 96.4 (1.0) & 96.4 (1.0) & 96.4 (1.0) & 96.4 (1.0) & 96.4 (1.0) \\
\bottomrule
\end{tabular}

\end{table}

\begin{table}[!htb]
  \centering\scriptsize
  \caption{E4 궤도 모드 학습(\code{mnist\_linear}, 각 10회). 학습 전체를 해당 일수의 LEO 노출로 가정. ``추론 모형 예측''은 추론 데이터로 맞춘 모델 A의 값.}
  \label{tab:e4leo}
  \begin{tabular}{rrr|rrr|rr}
\toprule
학습 시간 & 일수 & 기대 $k$ & \multicolumn{3}{c|}{방어 없음} & \multicolumn{2}{c}{클리핑} \\
(h) & & & 생존률 & 붕괴율 & 추론 모형 예측 & 생존률 & 붕괴율 \\
\midrule
24 & 1 & 3 & 1.0 & 0.1 & 0.97 & 1.0 & 0.0 \\
168 & 7 & 23 & 0.9 & 0.3 & 0.82 & 1.0 & 0.0 \\
720 & 30 & 98 & 0.4 & 0.9 & 0.43 & 1.0 & 0.0 \\
2160 & 90 & 293 & 0.2 & 1.0 & 0.08 & 1.0 & 0.0 \\
\bottomrule
\end{tabular}

\end{table}

%==================================================================
% 위치: 5 논의 — 한계 ‘수학 모델의 범위’ 항목
%==================================================================
학습 중 주입(E4)에서도 생존률은 같은 모형으로 예측되었지만, 실패의 모양(붕괴 대 부분 손상)과 클리핑의 효과는 모형 밖의 현상이다.

%==================================================================
% 위치: 6 결론 — 결과 (8)
%==================================================================
(8) 학습 도중의 명중은 치명 확률은 같되 실패가 전부 NaN 붕괴라 방어 없는 궤도상 학습은 수 일 안에 죽지만, 클리핑만 켜면 SGD가 손상을 고쳐 학습이 추론보다 강건해진다.

%==================================================================
% 위치: 6 향후 과제 — 궤도상 재학습 항목 (본문에는 일반화한 문장으로 대체)
%==================================================================
\item \textbf{궤도상 재학습}: E4에서 클리핑이 학습을 추론보다 강건하게 만든 결과를 큰 모델과 긴 학습(수십 epoch)에서 확인하고, 학습률과 클리핑 상수 $C$의 상호작용을 조사한다.

